% =====================================================================
%  Реферат
%  Оформление по ГОСТ 7.32-2017, п. 5.3, 6.12 и Приложение В.
%  Три компоненты: сведения об объеме, ключевые слова (прописными),
%  текст реферата с абзацного отступа.
% =====================================================================
\structelemnotoc{РЕФЕРАТ}

% Объем отчета указывается БЕЗ приложений (ГОСТ 7.32-2017, п. 6.12.1).
% Используется метка \label{lastpagebeforeappendix}, установленная
% в конце списка источников (файл 11_bibliography.tex).
Отчет \pageref{lastpagebeforeappendix}~с., 1~кн., 8~рис., 24~табл.,
37~источн., 3~прил.

\medskip

% Ключевые слова по ГОСТ 7.32-2017, п. 6.12.2:
% именительный падеж, прописными буквами, в строку, через запятые,
% без абзацного отступа и переноса слов, без точки в конце перечня.
\noindent
\begingroup
\hyphenpenalty=10000\exhyphenpenalty=10000\sloppy
ОБНАРУЖЕНИЕ ВТОРЖЕНИЙ, СЕТЕВАЯ БЕЗОПАСНОСТЬ, АНАЛИЗ СЕТЕВОГО
ТРАФИКА, ГЛУБОКОЕ ОБУЧЕНИЕ, НЕЙРОННЫЕ СЕТИ, CNN-LSTM, UNSW-NB15,
CICIDS2017, КОНЕЧНО-АВТОМАТНЫЙ АГЕНТ, ДРЕЙФ РАСПРЕДЕЛЕНИЙ,
ФОКАЛЬНАЯ ФУНКЦИЯ ПОТЕРЬ, ТЕМПЕРАТУРНОЕ МАСШТАБИРОВАНИЕ, PSI, MMD
\par
\endgroup

\medskip

Объектом исследования являются методы и технологии обнаружения
аномальной активности и несанкционированных действий в корпоративных
сетях на основе телеметрии сетевого трафика.

Цель работы --- разработка методики нейросетевого обнаружения
несанкционированных действий и аномальной активности в
корпоративной сети, валидируемой конечно-автоматным агентом
активного тестирования с библиотекой параметризованных шаблонов
атак.

В процессе работы выполнен систематический анализ литературы по
NIDS и глубокому обучению; проведена балльная оценка десяти
архитектур нейронных сетей по семи критериям и десяти публично
доступных датасетов NIDS по восьми критериям, опирающимся на
эталонный фреймворк Gharib и соавт. (2016) и обзор Ring и соавт.
(2019); реализована полная программная подсистема обработки данных,
обучения 15 моделей-кандидатов в едином регистре, инференса с
калибровкой порога, активного тестирования и мониторинга дрейфа;
проведено шесть экспериментов с 30 запусками и 15 моделями;
выполнен статистический анализ результатов через bootstrap-доверительные
интервалы 95\,\% и парный критерий Уилкоксона.

В результате исследования сформирована методика обнаружения,
включающая препроцессинг flow-данных, библиотеку из семи
параметризованных шаблонов атак, конечный автомат планирования
режимов норма/атака/дрейф (11 состояний) и блок drift-monitoring
на основе индексов PSI, KL и MMD. Реализован программный прототип
на стеке Python\,/\,PyTorch с интерфейсами FastAPI и Streamlit
(порядка 4000~строк исходного кода, 11 модулей юнит-тестов, 86
проходящих тестов). На датасете UNSW-NB15 проведены шесть
экспериментов (E1--E6) и cross-evaluation на CICIDS2017. По
композитному критерию качества целевая архитектура CNN-LSTM
лидирует по семи из десяти основных метрик: F1\,=\,0,9735,
PR-AUC\,=\,0,9945, ROC-AUC\,=\,0,9923, MCC\,=\,0,918,
Recall\,=\,0,985 на F1-max пороге; на операционном пороге под
целевую частоту ложных срабатываний 0,01 модель достигает
F1\,=\,0,971 при FPR\,=\,0,0095 и Recall\,=\,0,948. Inference-задержка
CNN-LSTM на CPU составила 1,16~мс/окно, пропускная способность ---
24\,560~EPS, что значительно превышает требования к
эксплуатационной производительности. Cross-evaluation
UNSW-NB15\,$\to$\,CICIDS2017 показал
$\Delta F_1$\,=\,0,18, что удовлетворяет требованию переносимости.
Все десять нефункциональных требований к системе выполнены.
Научная новизна состоит в интеграции нейросетевого детектора,
конечно-автоматного агента и блока мониторинга дрейфа в единую
воспроизводимую методику с прозрачной ground-truth по состояниям
FSM и единым протоколом экспериментов.

Областью применения результатов являются системы обнаружения
вторжений корпоративного класса, учебно-исследовательские стенды
по сетевой безопасности, задачи аудита устойчивости NIDS к дрейфу
распределений и вариативности атак.

Программный прототип допускает развертывание в составе подсистемы
NTA-аналитики корпоративного SOC. В отчете сформированы инженерные
рекомендации по интеграции с источниками NetFlow\,/\,IPFIX и
калибровке порогов под целевой эксплуатационный FPR. По сравнению
с лидером по F1 среди классических подходов (XGBoost, F1\,=\,0,965)
целевая модель обеспечивает в одиннадцать раз меньшее число ложных
срабатываний на типовом потоке корпоративного SOC за счет
существенно более низкого FPR на операционном пороге.

Эффективность работы определяется снижением частоты ложных
срабатываний детектора при сохранении полноты обнаружения за счет
температурной калибровки вероятностей и порога под допустимый FPR,
а также сокращением времени аналитика SOC на разбор инцидентов за
счет пятиступенчатой шкалы severity и временной дедупликации
алертов. Все поставленные задачи выпускной квалификационной работы
выполнены в полном объеме, гипотеза подтверждена, цель достигнута,
методика обнаружения сформулирована, программно реализована и
экспериментально проверена.
